\begin{frame}{실습 (3/3): 평가 $\varepsilon$의 비대칭}
  \begin{columns}[T]
    \begin{column}{0.58\textwidth}
      \centering
      \includegraphics[width=\linewidth]{ch08_2_eval_epsilon_asymmetry.png}
      \vskip0.3em
      {\footnotesize $\varepsilon=0\sim0.3$으로 200 에피소드 평가.
      $\varepsilon=0$이면 Q 앞선다($-13$ vs $-17$)가, 탐색이 조금만
      섞이면 Q 급락($\varepsilon=0.1\to -45$, $0.3\to -182$)해 SARSA가
      뒤집어 앞선다. SARSA 우회 경로는 거의 안 흔들림($-17\to -39$).}
    \end{column}
    \begin{column}{0.4\textwidth}
      \begin{block}{핵심 태도}
        \textbf{리뷰어는 발표팀의 실험을 \emph{자기 손으로} 재실행할
        수 있는 사람.} ``이 학습 곡선, 시드를 바꿔서 다시 돌려도
        같은 모양이 나올까요?''라는 질문은 \textbf{리뷰어가 실제로
        다시 돌려본 적 있는} 질문일 때 가장 날카롭다.
      \end{block}
      \vskip0.4em
      {\footnotesize ``평가 때 $\varepsilon$을 켜둔 채'' 비교하면
      \textbf{결론의 방향이 뒤집히는} 구조.}
    \end{column}
  \end{columns}
\end{frame}
